\subsection{Part A: Test Artifacts and Functional Testing}

\indent 

A simple online shopping feature is Login $\rightarrow$ Add to Cart $\rightarrow$ Checkout $\rightarrow$ Payment.

Based on this feature:

\bigskip

\noindent \textbf{Write 2 test cases for login functionality (Remember the components of a test case)}.

Example:

Test Case ID: TC\_Login\_01

Description: Verify that a registered user can log in successfully.

Preconditions: User must be registered and have valid credentials.

Test Steps:

Open the login page. Enter valid username and password. Click the “Login” button. Expected Result: User is redirected to the homepage/dashboard. Actual Result: To be filled during testing.

Status: Pass/Fail

Remarks: N/A

\bigskip

\noindent\textbf{Identify test data for valid, invalid, and boundary scenarios.}

Example:

Valid: username = user1@example.com, password = Password123

Invalid: username = user1@example.com, password = WrongPass

Boundary: Password exactly at minimum length (e.g., 8 characters), maximum length (e.g., 20 characters)

\bigskip

\noindent\textbf{Map each test case to the corresponding testing level (Unit, Integration, System, UAT).}

Example:

Unit Testing (for login function logic) and System Testing (end-to-end login validation)

\subsection{Part B: Performance & Non-Functional Testing Simulation}

\indent 

Scenario: E-commerce platform experiences slow checkout during Black Friday.

Based on this information:

\bigskip

\noindent\textbf{Which type of performance testing should be performed? (Load, Stress, Soak)}

Load Testing: To simulate normal and peak traffic (expected high number of concurrent users).

Stress Testing: To determine system limits under extreme traffic beyond expected peaks.

Soak Testing: Optional, to check system stability over long periods of high traffic.

Primary focus here: Load Testing (peak load simulation) and optionally Stress Testing.

\bigskip

\noindent\textbf{How would QA design tests to simulate 1000+ concurrent users?}

Example:

Use a performance testing tool (e.g., JMeter, LoadRunner, Gatling).

Define virtual users (VUs) to simulate 1000+ simultaneous users.

Create realistic user scenarios: login $\rightarrow$ browse $\rightarrow$ add to cart $\rightarrow$ checkout $\rightarrow$ payment.

Gradually ramp up users to identify system bottlenecks.

Run multiple iterations to mimic peak traffic conditions.

\bigskip

\noindent\textbf{What metrics would you monitor? (Response time, throughput, error rate)}

Example:

Response Time: Time to complete transactions or page load time.

Throughput: Number of transactions per second.

Error Rate: Failed transactions or failed server responses.

Resource Utilization: CPU, memory, network, database performance.

Latency: Delay in server response under load.

\subsection{Part C: Professional Issues \& QA Standards}

\indent

Scenario: A critical bug in a banking app is discovered hours before release. Management pressures the QA team to ignore it.

Based on this information:

\bigskip

\noindent\textbf{Identify ethical issues.}

Releasing software with a known critical bug could harm users or financial systems. QA professionals face pressure to compromise integrity for deadlines. Potential violations of professional ethics and legal compliance.

\bigskip

\noindent\textbf{Discuss the role of ISO 9000, ISO 27001, and QA governance in handling the situation.}

ISO 9000: Ensures a quality management system; mandates that defects are documented, addressed, and reviewed.

ISO 27001: Addresses information security, ensuring critical security flaws are fixed before release.

QA Governance: Provides structured risk-based decision-making, audits, and accountability, ensuring QA can escalate critical issues.

\bigskip

\noindent\textbf{Propose actions QA should take.}

Document the defect clearly, including impact assessment.

Escalate to management and stakeholders via formal channels.

Refuse to approve release until critical defect is resolved.

Recommend patch scheduling or release delay to ensure user safety and regulatory compliance.

Follow QA governance frameworks and ethical standards to maintain professional integrity.